% Part: lambda-calculus
% Chapter: introduction
% Section: computable-lambda

\documentclass[../../../include/open-logic-section]{subfiles}

\begin{document}

\olfileid[hi]{lam}{int}{lrp}
\olsection{संगणनीय फलन \usetoken{S}{lambda definable} हैं}

\begin{thm}
\ollabel{thm:computable-lambda}
प्रत्येक संगणनीय आंशिक फलन !!{lambda definable} है।
\end{thm}

\begin{proof}
हमें दिखाना है कि प्रत्येक आंशिक संगणनीय फलन~$f$, किसी लैम्ब्डा पद~$F$
द्वारा !!{lambda defined} है। क्लिनी के प्रसामान्य-रूप प्रमेय से यह दिखाना
पर्याप्त है कि प्रत्येक आदिम पुनरावर्ती फलन किसी लैम्ब्डा पद द्वारा
!!{lambda defined} है, और फिर यह कि !!{lambda definable} फलन उपयुक्त
संयोजनों और अपरिबद्ध खोज के अधीन संवृत हैं। प्रत्येक आदिम पुनरावर्ती फलन
को किसी लैम्ब्डा पद द्वारा !!{lambda defined} दिखाने के लिए यह दिखाना
पर्याप्त है कि प्रारंभिक फलन !!{lambda definable} हैं और !!{lambda
definable} आंशिक फलन संयोजन, आदिम पुनरावर्तन और अपरिबद्ध खोज के अधीन संवृत
हैं।
\end{proof}

शेष प्रमाण को अधिक पठनीय बनाने के लिए हम अधिक प्रचलित संकेत-लिपि उपयोग
करेंगे। उदाहरणार्थ, $Mxyz$ के बदले $M(x, y, z)$ लिखेंगे। यद्यपि यह
अर्थ-सूचक है, आपको याद रखना चाहिए कि टाइपरहित लैम्ब्डा कलन के पदों से कोई
नियत स्थानिकता संबद्ध नहीं होती; इसलिए उसी पद~$M$ के लिए $M(x,y)$ और
$M(x,y,z,w)$ लिखना भी उतना ही सार्थक है। किंतु यह संकेत-लिपि बताती है कि
हम $M$ को तीन चरों का फलन मान रहे हैं, और परिभाषाओं के पीछे का आशय अधिक
स्पष्ट करती है। इसी प्रकार, ``$M$ को $M =
\lambd[x][\lambd[y][\lambd[z][\dots]]]$ द्वारा परिभाषित करें'' के बदले हम
कहेंगे, ``$M$ को $M(x,y,z) = \dots$ द्वारा परिभाषित करें।''

\end{document}
